% ===================================================================
%  TABLO No 7 — Kışın köy. — İlkbaharda çiftlik.
%  Traduction 2026 d'apres texto/io, controlee sur texto/fr et
%  texto/hi. Consigne : voir texto/tr/10-tablo-01.tex.
% ===================================================================

%%K t07-noto-1 noto
\VUnotes{143.2mm}{%
(*) Yemeğin ayrıntısı için 6 ve 13 numaralı tablolara bakınız.%
}

%%K t07-apar-1 apar
\VUsaut{4.0mm}
\VUcentre{13.2pt}{90}{İKİNCİ DİZİ}
\VUsaut{3.0mm}
\VUfilet{16mm}
\VUsaut{5.6mm}
\VUtitre{15.0pt}{60}{TABLO No 7}
\VUsaut{3.0mm}

%%K t07-tit-1 sub
\VUcentre{12.0pt}{0}{\textit{Birinci sahne.}}
\VUsaut{3.0mm}

%%K t07-tit-2 sub
\VUpk{10.2pt}{40}{Kışın köy.}
\VUsaut{4.0mm}

%%K t07-c1-01-1 p
1. --- Yılbaşı tatilinde babamla birlikte Yeniköy'e gittim; babamın orada
görecek birkaç işi olan küçük bir köy.

%%K t07-c1-02-1 p
2. --- Şehirle o köy arasında işleyen \VUgras{posta arabasına}
\textsuperscript{(11)} bindik; ama hava çok soğuktu, ve böyle rahatsız bir
arabada yolculuk hiç de hoş olmadı.

%%K t07-c1-03-1 p
3. --- Bunun için \VUgras{arabacının} arabayı alanda \textit{Beyaz Ayı}
\VUgras{hanının} \textsuperscript{(2)} önünde durdurduğunu gördüğümüzde
gerçekten sevindik. \VUgras{Hancı} \textsuperscript{(4)} eşiğinde bizi
bekliyordu, sessizce \VUgras{piposunu} içerek.

%%K t07-c1-03-2 p
Bizi içeri çağırdı, ve ocakta çıtırdayan asma ateşinin önüne oturmam için bana
iki kez söylemesi gerekmedi. Sonra artık eski \VUgras{tabelayı}
\textsuperscript{(3)} gıcırdatan ve bacadan çıkan \VUgras{dumanı}
\textsuperscript{(1)} kara burgaçlar hâlinde savuran o titreten
\VUgras{poyrazı} hiç düşünmedim.

%%K t07-c1-04-1 p
4. --- Öğle yemeği vaktiydi. Ve gecikmeden, önümüze konan o yalın ama lezzetli
yemeğin hakkını tam verdik (*).

%%K t07-tit-3 sub
\VUpk{10.2pt}{40}{Birkaç ziyaret.}
\VUsaut{3.0mm}

%%K t07-c1-05-1 p
5. --- Yemekten sonra sigorta acentesi olan babamla birlikte müşterilerini
görmeye gittim. Önce hanın yanında oturan \VUgras{nalbanta}
\textsuperscript{(5)} gittik. Bir ata nal çakmakla uğraşıyordu. Yardımcısının
gücüne hayran oldum; nalbant ağır çekiç darbeleriyle \VUgras{nalı}
\textsuperscript{(6)} çakarken, o atın \VUgras{toynağını} deri önlüğünün
üstünde sımsıkı tutuyordu.

%%K t07-c1-06-1 p
6. --- Sonra köyün \VUgras{ayakkabı tamircisi} \textsuperscript{(7)} yaşlı
Yakup'a gittik. Tabelasında görkemli bir \VUgras{çizme} çizili olduğu hâlde,
eski yıpranmış ayakkabılara yeni pençe vurarak geçiniyordu.

%%K t07-c1-06-2 p
Yaşlı adam gülerek dedi ki: « Şehirde “ayakkabıcı” idim ve hiç müşterim
yoktu. Burada “tamirciyim” ve iş hiç eksilmiyor. »

%%K t07-c1-07-1 p
7. --- Bize komşusu \VUgras{berberden} \textsuperscript{(8)} söz etti, ki bütün
müşterileri hoşnutsuz. \VUgras{Usturaları} hep kertiklidir ve
\VUgras{makası} hiçbir zaman doğru kesmez. Ama köyün tek berberi o olduğundan,
insanların ona tıraş olmaktan ve saç kestirmekten başka çaresi yok. Üstelik
cimri ve sert.

%%K t07-c1-08-1 p
8. --- Bereket öteki \VUgras{komşular} onun gibi değil. Söz gelişi
\VUgras{arabacı} \textsuperscript{(9)} iyi bir adam, çalışkan ve
yardımsever. Arabayı ona onarmaya vermek bir zevk. İşi her zaman temiz.

%%K t07-c1-09-1 p
9. --- Yaşlı Yakup'un \VUgras{saraçtan} \textsuperscript{(10)} da bir yakınması
yoktu, ki işin sıkıştığı zamanlarda ona \VUgras{koşum takımı} dikmekte yardım
eder.

%%K t07-c1-09-2 p
Babam ona ailesini sordu. « Hepsi sağ salim. Torunum \VUgras{okulda}
\textsuperscript{(12)}. \VUgras{Öğretmeni} \textsuperscript{(13)} ondan pek
hoşnut; iyi bir \VUgras{öğrenci} \textsuperscript{(14)}. Şu yaramaz
oğlum gibi değil; oyundan başka bir şey düşünmüyor. \VUgras{Baş öğretmen}
\textsuperscript{(15)} ona hiçbir şey öğretemiyor. »

%%K t07-tit-4 sub
\VUsaut{3.0mm}
\VUpk{10.2pt}{40}{Köy haberleri.}
\VUsaut{3.0mm}

%%K t07-c1-10-1 p
10. --- Sonra yaşlı adam bize köyün haberlerini anlattı. Yeni bir okul
yapılacaktı, çünkü \VUgras{köy konağındaki} çok küçük gelmişti. Bu kolay
olurdu, çünkü \VUgras{değirmencinin} bahçesinin bir bölümünü satın almak
yeterdi. Zavallı için bu çok iyi olurdu. Soğuklar bastıralı, \VUgras{değirmenin}
\textsuperscript{(16)} \VUgras{taşları} buğday öğütemedi, çünkü
\VUgras{çark} \textsuperscript{(17)} \VUgras{buza} \textsuperscript{(25)}
saplanmıştı, ki \VUgras{derenindi} \textsuperscript{(23)}.

%%K t07-c1-11-1 p
11. --- « Yapı demişken --- yeni \VUgras{kilisemizi} \textsuperscript{(18)}
gördünüz mü? Kar örtüsünün altında çok görünüşlü durur.
\VUgras{Kargalar} \textsuperscript{(19)}, ki eski çan kulelerini artık
tanımıyorlar, üzgün üzgün \VUgras{mezarlığın} \textsuperscript{(20)} büyük
ağacına konup duruyorlar. »

%%K t07-c1-12-1 p
12. --- Mezarlığın adı geçince tamirci, babamın tanıdığı ve geçen yıl göçüp
gitmiş kimseleri anımsadı. « Ne çok \VUgras{mezar} \textsuperscript{(21)},
sevgili beyefendi, \VUgras{servilerin} \textsuperscript{(22)} altında eklendi!
Ölüm yaşlı kâtibimizi aldı. Bütün köy \VUgras{cenazesindeydi}. »

%%K t07-c1-13-1 p
13. --- « İşte onun ardılı, derede buz üstünde kayıyor. Biraz önce gelmişti ki
kopan \VUgras{buz patenlerinin} \textsuperscript{(26)} bağını ekleyeyim. »

%%K t07-c1-14-1 p
14. --- « Ve işte, derenin kıyısında, yeni \VUgras{köy bekçisi}
\textsuperscript{(27)}. Eski bir askerdir ve köyün yaramaz oğlanlarının
belası. \VUgras{Kayışı} \textsuperscript{(29)} ile \VUgras{püsküllü
kepi} \textsuperscript{(28)} görünür görünmez kaçışırlar. »

%%K t07-c1-15-1 p
15. --- « Vay! İşte yaşlı Yusuf, fırıncıya \VUgras{odun demeti arabasını}
\textsuperscript{(30)} getiriyor. --- Doğru, dedi babam, \VUgras{katır}
\textsuperscript{(31)} çiftini tanıyorum. Onunla konuşacak bir işim var, tam
fırsat. Hoşça kal, yaşlı Yakup. »

%%K t07-c1-16-1 p
16. --- Tam çıkıyorduk ki köyün çocukları okuldan salınıp \VUgras{buz
kaydırağına} \textsuperscript{(33)} üşüştüler, \VUgras{düşüp}
\textsuperscript{(34)} kollarını bacaklarını kırma tehlikesiyle. İçlerinden
biri arabacıdan \VUgras{kızağını} \textsuperscript{(32)} getirdi, üstüne
arkadaşlarından birini oturttu, ve koşmaya başladı.

%%K t07-c1-17-1 p
17. --- Ötekiler bir \VUgras{kar yığınına} \textsuperscript{(37)} atıldılar, ki
\VUgras{köprünün} \textsuperscript{(40)} yanındaydı, ve dün yaptıkları
\VUgras{kardan adama} \textsuperscript{(39)} vurmaya koyuldular; ona bir
şapka, bir pipo ve omzuna asılı bir çanta vermişlerdi.

%%K t07-c1-18-1 p
18. --- Buzlu yeli ve soğuğu umursamıyorlardı. Çoğunun bir \VUgras{atkısı}
\textsuperscript{(35)} bile yoktu, ve \VUgras{kar topu}
\textsuperscript{(38)} savaşından sonra ısınmak için bir avuç sıcacık
\VUgras{kestane} \textsuperscript{(42)} onlara yetiyordu; onu da yaşlı Jaklin
\VUgras{satıcı kadından} alıyorlardı, ki çok ince \VUgras{şalı}
\textsuperscript{(43)} içinde soğuktan titrer durur.

%%K t07-tit-5 sub
\VUsaut{3.0mm}
\VUcentre{12.0pt}{0}{\textit{İkinci sahne.}}
\VUsaut{3.0mm}

%%K t07-tit-6 sub
\VUpk{10.2pt}{40}{İlkbaharda çiftlik evi.}
\VUsaut{4.0mm}

%%K t07-c2-01-1 p
1. --- Kışın bütün tatlarına karşın \VUgras{ilkbaharın} dönüşünü derin bir
sevinçle karşılarız.

%%K t07-c2-01-2 p
Mayıs ayında akşamüstü çiftliğin avlusu ne coşkulu bir tablo yapar!

%%K t07-c2-02-1 p
2. --- Ufukta \VUgras{güneş} \textsuperscript{(1)} ağır ağır batıyor, ve
\VUgras{kırları} \textsuperscript{(3)} kızıl \VUgras{ışınlarıyla}
\textsuperscript{(2)} dolduruyor. Çalışkan \VUgras{çiftçi}
\textsuperscript{(6)} \VUgras{tarlasını} \textsuperscript{(4)} sürmeye
acele ediyor.

%%K t07-c2-02-2 p
\VUgras{Sabanıyla} \textsuperscript{(5)}, ki güçlü bir \VUgras{at}
\textsuperscript{(7)} çeker, \VUgras{karıkları} \textsuperscript{(8)}
çiziyor;
\VUgras{ekici} \textsuperscript{(9)} onlara avuç avuç o \VUgras{tohumu}
serpecek, ki ondan altın başaklar çıkacak.

%%K t07-c2-03-1 p
3. --- Tırpanlarıyla \VUgras{ot biçenler} \textsuperscript{(11)} çayırın
otunu biçmiş, kurutucular \VUgras{kuru otu} \textsuperscript{(10)}
\VUgras{tırmıklarla} \textsuperscript{(12)} ve dirgenlerle altüst edip
kurutmuş, ve şimdi dönüyorlar.

%%K t07-c2-03-2 p
Kimi öküzlerin çektiği ağır arabanın önünde yürüyor, omzunda aletiyle.
Ötekiler, daha tembel, kokulu kuru otun üstüne rahatça serilmiş.

%%K t07-c2-04-1 p
4. --- Geçerken \VUgras{bahçıvanı} \textsuperscript{(16)} dostça selamlıyorlar,
ki \VUgras{bahçede} \textsuperscript{(13)} uğraşır, \VUgras{meyve ağaçlarını}
budar ve \VUgras{armut ağaçlarına} \textsuperscript{(14)} aşı yapar.
\VUgras{Erik ağaçları} \textsuperscript{(15)} çiçek yüklü, ve iyi gübrelenmiş
\VUgras{sebze bahçesinde} \textsuperscript{(17)} sebzeler, lahanalar ve
salatalar şimdiden güzel oluyor.

%%K t07-c2-04-2 p
Alçak duvarın üstünde güzel bir \VUgras{tavus kuşu} \textsuperscript{(18)}
görkemli tüylerini göstermek için kuyruğunu açıyor.

%%K t07-tit-7 sub
\VUpk{10.2pt}{40}{Çiftliğin avlusu.}
\VUsaut{3.0mm}

%%K t07-c2-05-1 p
5. --- \VUgras{Ahırın} \textsuperscript{(19)} kapıları açık, ve yemlik
ile \VUgras{yataklık saman} \textsuperscript{(23)} görünür, ki
\VUgras{kısrağın} \textsuperscript{(21)}; kısrağı \VUgras{seyis}
\textsuperscript{(20)} az önce çözmüş. \VUgras{Tay} \textsuperscript{(22)}, uzun bacakları
üstünde öyle tuhaf, sıçraya sıçraya annesinin ardından gidiyor.

%%K t07-c2-06-1 p
6. --- \VUgras{Kuyunun} \textsuperscript{(29)} yanında \VUgras{inekler}
\textsuperscript{(25)} tahtadan \VUgras{yalaktan} \textsuperscript{(26)} su
içmeye gidiyor, ki onların çeşmesi. \VUgras{Buzağı}
\textsuperscript{(24)} bu fırsatta sütünü içiyor, ve \VUgras{öküz}
\textsuperscript{(27)}, yemin güzel kokusunu alıp, sevinçle böğürüyor ve güçlü
\VUgras{boynuzlu} \textsuperscript{(28)} başını gururla kaldırıyor.

%%K t07-c2-07-1 p
7. --- Herkes iş başında. \VUgras{Hizmetçi kız} \textsuperscript{(32)}
kuyunun \VUgras{ipini} \textsuperscript{(31)} tutuyor. Hayvanları sulamak için
bir \VUgras{kovaya} \textsuperscript{(30)} su çekiyor. \VUgras{Samanlığın}
\textsuperscript{(33)} yanında bir adam dağılmış samanı topluyor, ve
\VUgras{çiftlik evinin} \textsuperscript{(34)} kapısı önünde yaşlı bir
\VUgras{eğirici kadın} \textsuperscript{(35)}, çıkrığı ve \VUgras{iğiyle},
eski günlerdeki gibi \VUgras{keten} ya da \VUgras{kenevir} eğiriyor.

%%K t07-tit-8 sub
\VUsaut{3.0mm}
\VUpk{10.2pt}{40}{Çiftliğin sahibi.}
\VUsaut{3.0mm}

%%K t07-c2-08-1 p
8. --- \VUgras{Çiftliğin sahibi} \textsuperscript{(36)} bütün bu işi yürütür ve
adamlarına buyruk verir. Onlara \VUgras{tapanı} \textsuperscript{(40)} ile
\VUgras{çapayı} \textsuperscript{(39)} örtüp kaldırmalarını söylüyor; onlar da
\VUgras{kulübenin} \textsuperscript{(38)} yanında bırakılmış, ki
\VUgras{köpeğin} \textsuperscript{(37)}.

%%K t07-c2-09-1 p
9. --- Sesiyle bir \VUgras{güvercin} \textsuperscript{(41)} ürküyor, ki tam
hızla \VUgras{güvercinliğe} \textsuperscript{(42)} uçuyor, ve
\VUgras{tavuklar} \textsuperscript{(43)} da, ki telaşla \VUgras{kümese}
\textsuperscript{(44)} dönüyor.

%%K t07-c2-10-1 p
10. --- \VUgras{Ağılın} \textsuperscript{(48)} önünde, işte yaşlı
\VUgras{çoban} \textsuperscript{(45)} \VUgras{sürüsünü} \textsuperscript{(49)}
döndürüyor. \VUgras{Değneğini} \textsuperscript{(46)} tutarak, ki solmuş
\VUgras{gömleği} \textsuperscript{(47)} kadar yanından ayrılmaz, ağıla doğru
\VUgras{koçları} \textsuperscript{(50)}, \VUgras{koyunları}
\textsuperscript{(53)}, \VUgras{kuzuları} \textsuperscript{(54)},
\VUgras{keçileri} \textsuperscript{(51)} ve \VUgras{oğlakları}
\textsuperscript{(52)} sürüyor. Sadık \VUgras{köpeği} \textsuperscript{(55)}
Argus en arkadan geliyor ve havlayarak geri kalanları öne katıyor.

%%K t07-c2-11-1 p
11. --- Bütün bu gürültü patırtı o iyi \VUgras{sağmal ineğin}
\textsuperscript{(56)} rahatını bozmaz, ki \VUgras{çanını}
\textsuperscript{(57)} çalıp durur; o inek \VUgras{ineklerin ahırının}
\textsuperscript{(58)} kapısı önünde sağılır.

%%K t07-c2-12-1 p
12. --- Sanki anlıyor ki sütünü vermesi gerekir, \VUgras{çiftliğin sahibesi}
\textsuperscript{(60)} \VUgras{yayıkta} \textsuperscript{(61)}
\VUgras{tereyağı} yapabilsin, ya da o güzel \VUgras{peynirleri}
\textsuperscript{(64)}, ki \VUgras{teknenin} \textsuperscript{(63)} üstündeki
tahtada süzülür.

%%K t07-c2-13-1 p
13. --- Bütün \VUgras{süthane} \textsuperscript{(59)} çok temiz durur,
\VUgras{çiftlik işçisinin} \textsuperscript{(65)} sayesinde; o da az önce
taşıdığı büyük kovada \VUgras{süt güğümlerini} \textsuperscript{(62)}
yıkamış.

%%K t07-tit-9 sub
\VUsaut{3.0mm}
\VUpk{10.2pt}{40}{Kümes.}
\VUsaut{3.0mm}

%%K t07-c2-14-1 p
14. --- Ağır takunyalarıyla biraz hantal yürüyor, ve bundan
\VUgras{hindiler} \textsuperscript{(68)}, \VUgras{ördekler}
\textsuperscript{(66)}, \VUgras{suna ördekleri} \textsuperscript{(67)} ve
\VUgras{kazlar} \textsuperscript{(69)} kaçışıyor; \VUgras{gagalarını}
\textsuperscript{(70)} ardına kadar açıp tedirgin bir gürültü koparıyorlar.

%%K t07-c2-14-2 p
Daha kavgacı olan \VUgras{horoz} \textsuperscript{(71)}, kızıl
\VUgras{ibiğini} \textsuperscript{(72)} kaldırıyor, \VUgras{kuyruğunun}
\textsuperscript{(73)} tüylerini silkeliyor ve çınlayan bir « ü-ürü-üü »
salıveriyor.

%%K t07-c2-15-1 p
15. --- Bir \VUgras{ana tavuk} \textsuperscript{(74)}
\VUgras{gübre yığınında} \textsuperscript{(81)} eşiniyor, \VUgras{civcivleriyle}
\textsuperscript{(75)} birlikte. \VUgras{Domuz yavrusu}
\textsuperscript{(78)} annesine doğru koşuyor, ağılın yanında gördüğümüz o iri
\VUgras{dişi domuza} \textsuperscript{(77)} doğru.

%%K t07-c2-16-1 p
16. --- Kafeslerinde \VUgras{tavşanlar} \textsuperscript{(79)} sahibin kızının
kendilerine getirdiği o taze \VUgras{otu} \textsuperscript{(80)} iştahla
yiyorlar. Jeni tavşanlarını çok sever, ve her gün, kümesten taze
\VUgras{yumurtaları} \textsuperscript{(76)} aldıktan sonra, küçük dostlarını
görmeye gelir.
\VUsaut{6.0mm}
\VUfilet{22mm}
